% Part: first-order-logic
% Chapter: models-theories
% Section: introduction

\documentclass[../../../include/open-logic-section]{subfiles}

\begin{document}

% SEGMENT OLP-0168-S01

\olfileid{fol}{mat}{int}
\olsection{அறிமுகம்}

\begin{explain}
அடிகோள் முறையின் வளர்ச்சி அறிவியல் வரலாற்றில் ஒரு குறிப்பிடத்தக்க
சாதனை; கணித வரலாற்றில் அது தனிச் சிறப்புடையது. ஒரு துறையின் அடிகோள்
வழி வளர்ச்சி பல கேள்விகளைத் தெளிவாக்குகிறது: அந்தத் துறை எதைப்
பற்றியது? மிக அடிப்படையான கருத்துகள் எவை? அவை எவ்வாறு தொடர்புடையவை?
அந்தத் துறையின் எல்லாக் கருத்துகளையும் இவ்வடிப்படைக் கருத்துகளைக் கொண்டு
வரையறுக்க முடியுமா? இக்கருத்துகள் எந்த விதிகளுக்குக் கீழ்ப்படிகின்றன,
எவற்றுக்குக் கீழ்ப்படிந்தே ஆக வேண்டும்?

அடிகோள் முறையும் தருக்கமும் ஒன்றுக்கொன்று இயல்பாகப் பொருந்துகின்றன.
அடிகோள் கோட்பாடுகளை உருவாக்கவும், கோட்பாட்டின் அடிகோள்களிலிருந்து
தேற்றங்களைத் துல்லியமாகக் குறிப்பிடப்பட்ட முறையில் நிரூபிக்கவும்,
அடிகோள்களை நிறைவுறுத்தும் எல்லா அமைப்புகளின் பண்புகளை முறையாக
ஆராயவும் முறையிய தருக்கம் கருவிகளை வழங்குகிறது.
\end{explain}

\begin{defn}
!!{sentence}s இன் தொகுப்பு~$\Gamma$ \emph{மூடியது} iff,
$\Gamma \Entails !A$ எனும்போதெல்லாம் $!A \in \Gamma$.
!!{sentence}s இன் தொகுப்பு~$\Gamma$ இன் \emph{மூடல்}
$\Setabs{!A}{\Gamma \Entails !A}$ ஆகும்.

$\Gamma$ ஆனது $\Delta$ ஆல் \emph{அடிகோளாக்கப்படுகிறது} என்கிறோம் iff,
$\Gamma$ ஆனது வாக்கியங்களின் தொகுப்பு~$\Delta$ இன் மூடலாக இருக்கும்.
\end{defn}

\begin{explain}
ஒரு அடிகோள் கோட்பாட்டை, அதன் அடிகோள்களின் தொகுப்பு~$\Delta$ ஆல்
அடிகோளாக்கப்படும் வாக்கியங்களின் தொகுப்பாகக் கருதலாம். அதாவது, நாம்
ஆராய விரும்பும் அடிகோள் வழி வளர்க்கப்பட்ட அறிவியலின் அடிப்படைக்
கருத்துகளுக்கான தருக்கமல்லாக் குறியீடுகளைக் கொண்ட ஒரு முதல்தர மொழியும்,
அந்த அறிவியலின் அடிப்படை விதிகளை வெளிப்படுத்தும் !!{sentence}s இன்
தொகுப்பும் நம்மிடம் இருக்கும்போது, அடிகோள்களால் உட்கிடையச் செய்யப்படும்
அந்த மொழியின் எல்லா !!{sentence}s உம் கோட்பாட்டைப் பிரதிநிதித்துவப்படுத்துகின்றன
என்று கருதலாம். பகுதி வரிசைகளின் கோட்பாடு போன்ற ஒரே அடிப்படைக் கருத்தும்
எளிய அடிகோள்களும் கொண்ட எடுத்துக்காட்டுகளிலிருந்து நியூட்டனிய இயந்திரவியல்
போன்ற சிக்கலான கோட்பாடுகள் வரை இது விரிகிறது.

அடிகோள் முறைக்கான இம்முறையிய அணுகுமுறையை முக்கியமாக்கும் தருக்க
உண்மைகள் பின்வருவன. $\Gamma$ என்பது ஒரு கோட்பாட்டுக்கான அடிகோள்
அமைப்பு, அதாவது வாக்கியங்களின் தொகுப்பு, என்க.
\begin{enumerate}
\item ஓர் அடிகோள் அமைப்பு, நோக்கமாகக் கொண்ட !!{structure}s இன் ஒரு
  வகுப்பை எப்போது சரியாகப் பிடிக்கிறது என்பதைத் துல்லியமாகக் கூறலாம்.
  அதாவது, !!{structure}s இன் ஒரு குறிப்பிட்ட வகுப்பில் நாம் ஆர்வம்
  கொண்டிருந்தால், $\Sat{M}{\Gamma}$ ஆகும்~$\Struct M$ களே அந்த
  !!{structure}s iff, அடிகோள் அமைப்பு~$\Gamma$ அவ்வகுப்பை வெற்றிகரமாகப்
  பிடிக்கிறது.
\item $\Sat{M}{\Gamma}$ ஆகும் $\Struct M$ ஒன்று இருந்தும், அந்த
  $\Struct M$ நாம் நோக்கமாகக் கொண்ட !!{structure}s இல் ஒன்றாக இல்லாததால்
  இவ்வகையில் நாம்
  தோல்வியடையலாம். அப்போது~$\Struct M$ இல் மெய்யாக இல்லாத அடிகோள்களைச்
  சேர்க்க நேரலாம்.
\item நோக்கமாகக் கொண்ட எல்லா !!{structure}s இலும் $\Gamma$ மெய் என்ற
  அளவிற்காவது நாம் வெற்றிபெற்றால், $\Gamma \Entails !A$ எனும்போதெல்லாம்
  வாக்கியம்~$!A$ நோக்கமாகக் கொண்ட எல்லா !!{structure}s இலும் மெய்யாகும்.
  எனவே, வாக்கியங்கள் அடிகோள்களால் உட்கிடையச் செய்யப்படுகின்றன என்று
  காட்டுவதன் மூலமே, தருக்கக் கருவிகளைப் (எடுத்துக்காட்டாக,
  !!{derivation} முறைகளைப்) பயன்படுத்தி அவை நோக்கமாகக் கொண்ட எல்லா
  !!{structure}s இலும் மெய்யானவை என்று காட்டலாம்.
\item சிலவேளைகளில் நோக்கமாகக் கொண்ட !!{structure}s எதையும் மனத்தில்
  கொள்ளாமல், அடிகோள்களிலிருந்தே தொடங்குகிறோம்: சில விதிகளை
  நிறைவுறுத்த வேண்டும் என்று விரும்பும் சில அடிப்படைக் கருத்துகளை
  எடுத்துக்கொண்டு, அவ்விதிகளை ஓர் அடிகோள் அமைப்பாகக் குறிமுறைப்படுத்துகிறோம்.
  அடிகோள்கள் ஒன்றோடொன்று முரண்படவில்லை என்பதை உடனே சரிபார்க்க விரும்புவோம்.
  அவை முரண்பட்டால், அவ்விதிகளுக்குக் கீழ்ப்படியும் கருத்துகள் எதுவும்
  இருக்க முடியாது; அப்போது நாம் ஒத்திசைவற்ற ஒரு கோட்பாட்டை அமைக்க
  முயன்றிருப்போம். $\Gamma$ இன் ஒரு மாதிரியைக் கண்டறிவதன் மூலம் இது
  நிகழவில்லை என்பதைச் சரிபார்க்கலாம். நமது கோட்பாட்டிற்கு மாதிரிகள்
  இருந்தால், தருக்க முறைகளால் அவற்றை ஆராயவும் புதிய மாதிரிகளை
  உருவாக்கவும் முடியும்.
\item அடிகோள்களின் சார்பின்மையும் ஒரு முக்கியமான கேள்வி. ஓர் அடிகோள்
  மற்றவற்றின் பின்விளைவாக இருந்து, எனவே தேவையற்றதாக இருக்கலாம்.
  $\Gamma$ இல் உள்ள அடிகோள் $!A$ தேவையற்றது என்பதை
  $\Gamma \setminus \{!A\} \Entails !A$ என்று நிரூபித்துக் காட்டலாம்.
  ஓர் அடிகோள் தேவையானது என்பதை $(\Gamma \setminus \{!A\}) \cup \{\lnot
  !A\}$ நிறைவுசெய்யத்தக்கது என்று காட்டியும் நிரூபிக்கலாம்.
  எடுத்துக்காட்டாக, வடிவியலின் மற்ற அடிகோள்களிலிருந்து இணைநேர்க்கோட்டு
  அடிகோள் சார்பற்றது என்பது இவ்வாறே காட்டப்பட்டது.
\item ஒரு கோட்பாட்டில் கருத்துகளை வரையறுக்க முடியுமா என்பதும் மற்றொரு
  முக்கியமான கேள்வி. மொழியின் தேர்வு, ஒரு கோட்பாட்டின் மாதிரிகளில்
  என்னென்ன உள்ளன என்பதைத் தீர்மானிக்கிறது. ஆனால் கோட்பாட்டின் ஒவ்வோர்
  அம்சமும் அதன் மாதிரிகளில் தனியாகப் பிரதிநிதித்துவப்படுத்தப்பட வேண்டியதில்லை.
  எடுத்துக்காட்டாக, ஒவ்வொரு வரிசை~$\le$ உம் அதற்குரிய கடு வரிசை~$<$
  ஐத் தீர்மானிக்கிறது---ஒன்று கொடுக்கப்பட்டால் மற்றதை வரையறுக்கலாம்.
  ஆகவே அத்தகைய வரிசையைக் கொண்ட கோட்பாட்டின் மாதிரி, அதற்குரிய கடு
  வரிசையையும் \emph{கூடுதலாகக்} கொண்டிருக்க வேண்டியதில்லை. பொதுவாக,
  ஒரு தொடர்பை மற்ற தொடர்புகளைக் கொண்டு எப்போது வரையறுக்க முடியும்?
  எப்போது ஒரு தொடர்பை மற்றவற்றைக் கொண்டு வரையறுக்க முடியாது
  (அதனால் அதை மொழியின் அடிப்படைக் குறியீடுகளில் சேர்க்க வேண்டும்)?
\end{enumerate}
\end{explain}


\end{document}
